\chapter{Travail réaliser}

Ce chapitre se concentrera sur le travail réaliser durant ce travail de Bachelor.

\section{Objectif à réaliser}

Les différents objectifs à atteindre sont les suivants :

\begin{itemize}
    \item Trouver le problème du code principal
    \item Réaliser l'interface utilisateur
    \item Réaliser une synthèse d'un régulateur PID
    \item
\end{itemize}

\section{Recherche du problème lié au code principal}

\subsection{Hypothèse}

Les hypothèses suivantes ont été pensée qui causerait potentiellement le problème :

\begin{itemize}
    \item Les interruptions liées à la communication sans fil empêchent le bon déroulement de l'interruption de la régulation.
\end{itemize}

Les sous-sections suivantes viennent traiter les hypothèses établies.

\subsection{Vérifications des interruptions}

Afin de vérifier si les interruptiosn liée à la communication sans fil viendrait à gêner l'interruption principal (ADC, régulation, PWM,...), il suffit simplement de mettre en commentaire les parties concernant la communication sans fil et flashé le DSP afin de vérifier l'hypothèses.\\

Les parties mises en commentaire sont les suivantes :

\begin{minted}{c}
    __interrupt void INT_mySCI0_RX_ISR(void) // .xxx
    {
        uint16_t c;

        while (SCI_getRxFIFOStatus(SCIA_BASE) > 0) {
            c = SCI_readCharBlockingFIFO(SCIA_BASE);
            if (rxIndex < RX_BUF_LEN) {
                rxBuffer[rxIndex++] = c;
            }
            if(c == '\x02'){
                // Marque le début de trame
                dataIndex = rxIndex; // Sauvegarde la position pour les données
            }
            if(c == '\x03'){
                // Fin de trame - traite les données
                if(dataIndex > 0 && dataIndex < rxIndex-1) {
                    state_PIN = (rxBuffer[dataIndex] == '1') ? 1 : 0;
                }
                // Reset du buffer
                rxIndex = 0;
                int b;
                for(b = 0; b < RX_BUF_LEN; b++) {
                    rxBuffer[b] = 0;
                }
            }
        }

        SCI_clearOverflowStatus(SCIA_BASE);
        SCI_clearInterruptStatus(SCIA_BASE, SCI_INT_RXFF);
        Interrupt_clearACKGroup(INTERRUPT_ACK_GROUP9);
    }

    __interrupt void INT_mySCI0_TX_ISR(void)
    {
        while (SCI_getTxFIFOStatus(SCIA_BASE) < SCI_FIFO_TX16 && txIndex < txLength) {
            SCI_writeCharBlockingFIFO(SCIA_BASE, txBuffer[txIndex++]);
        }

        if (txIndex >= txLength) {
            SCI_disableInterrupt(mySCI0_BASE, SCI_INT_TXFF); // Fin d’envoi
        }

        // Acquitter l'interruption
        SCI_clearOverflowStatus(SCIA_BASE);
        SCI_clearInterruptStatus(SCIA_BASE, SCI_INT_TXFF);
        Interrupt_clearACKGroup(INTERRUPT_ACK_GROUP9); // For the PIE acknowledge
    }
\end{minted}

\section{Interface utilisateur}

\section{Synthèse du régulateur PID}
